\icode{N23L-MC-SI-L2-261026-003}
\nceeID[4]{2023辽宁}{1}
某同学在练习投篮，篮球在空中的运动轨迹如图中虚线所示，篮球所受合力\(F\)的示意图可能正确的是\paren[]

\begin{tikzpicture}[>=stealth,font=\scriptsize]
	\begin{scope}
		\draw[name path=curve,dashed] (3,2) parabola bend(3,2) (0,0);
		\path[name path=beeline] (1.6,0)--++(0,2);
		\path[name intersections={of=curve and beeline,by=p}];
		\draw[brown] (p) circle(.25);
		\draw[brown] ($(p)+(0,-.25)$) arc(270:90:.075 and .25) ($(p)+(-.25,0)$) arc(-180:0:.25 and .075);
		\draw[brown] ($(p)+(135:.25)$)..controls($(p)+(-.125,-.125)$)and($(p)+(.125,.25)$)..($(p)+(45:.25)$);
		\draw[brown] ($(p)+(-135:.25)$)..controls($(p)+(-.125,.06)$)and($(p)+(.125,-.25)$)..($(p)+(-45:.25)$);
		\draw[dashed,cyan] (p)--++(0,-1);
		\draw[->,magenta] (p)--++(-110:.9)node[below]{ \(F\)};
		\node at(1.5,-.3){\normalsize A};
	\end{scope}
	\begin{scope}[xshift=3.5cm]
		\draw[name path=curve,dashed] (3,2) parabola bend(3,2) (0,0);
		\path[name path=beeline] (1.6,0)--++(0,2);
		\path[name intersections={of=curve and beeline,by=p}];
		\draw[brown] (p) circle(.25);
		\draw[brown] ($(p)+(0,-.25)$) arc(270:90:.075 and .25) ($(p)+(-.25,0)$) arc(-180:0:.25 and .075);
		\draw[brown] ($(p)+(135:.25)$)..controls($(p)+(-.125,-.125)$)and($(p)+(.125,.25)$)..($(p)+(45:.25)$);
		\draw[brown] ($(p)+(-135:.25)$)..controls($(p)+(-.125,.06)$)and($(p)+(.125,-.25)$)..($(p)+(-45:.25)$);
		\draw[dashed,cyan] (p)--++(0,-1);
		\draw[->,magenta] (p)--++(90:.9)node[right]{ \(F\)};
		\node at(1.5,-.3){\normalsize B};
	\end{scope}
	\begin{scope}[xshift=7cm]
		\draw[name path=curve,dashed] (3,2) parabola bend(3,2) (0,0);
		\path[name path=beeline] (1.6,0)--++(0,2);
		\path[name intersections={of=curve and beeline,by=p}];
		\draw[brown] (p) circle(.25);
		\draw[brown] ($(p)+(0,-.25)$) arc(270:90:.075 and .25) ($(p)+(-.25,0)$) arc(-180:0:.25 and .075);
		\draw[brown] ($(p)+(135:.25)$)..controls($(p)+(-.125,-.125)$)and($(p)+(.125,.25)$)..($(p)+(45:.25)$);
		\draw[brown] ($(p)+(-135:.25)$)..controls($(p)+(-.125,.06)$)and($(p)+(.125,-.25)$)..($(p)+(-45:.25)$);
		\draw[dashed,cyan] (p)--++(0,-1);
		\draw[->,magenta] (p)--++(120:.9)node[left]{ \(F\)};
		\node at(1.5,-.3){\normalsize C};
	\end{scope}
	\begin{scope}[xshift=10.5cm]
		\draw[name path=curve,dashed] (3,2) parabola bend(3,2) (0,0);
		\path[name path=beeline] (1.6,0)--++(0,2);
		\path[name intersections={of=curve and beeline,by=p}];
		\draw[brown] (p) circle(.25);
		\draw[brown] ($(p)+(0,-.25)$) arc(270:90:.075 and .25) ($(p)+(-.25,0)$) arc(-180:0:.25 and .075);
		\draw[brown] ($(p)+(135:.25)$)..controls($(p)+(-.125,-.125)$)and($(p)+(.125,.25)$)..($(p)+(45:.25)$);
		\draw[brown] ($(p)+(-135:.25)$)..controls($(p)+(-.125,.06)$)and($(p)+(.125,-.25)$)..($(p)+(-45:.25)$);
		\draw[dashed,cyan] (p)--++(0,-1);
		\draw[->,magenta] (p)--++(-155:.9)node[left]{ \(F\)};
		\node at(1.5,-.3){\normalsize D};
	\end{scope}
\end{tikzpicture}
\begin{answer}
	A\qquad 提示：篮球在空气中受空气阻力
\end{answer}